\chapter{Convexity of the Incomplete-Information Value}
\label{sec:p17-astrom-incomplete-information-ii}

\AA str\"om's 1969 Part II proves that the finite-horizon
maximum-reward value is convex in the information state;
the minimum-cost value is concave
\cite[pp.~403--406]{Astrom1969IncompleteInformationII}.
It uses positive linear observation operators and a change
of weights in a normalized convex combination.
This is a separate presentation from the filtering and
control theorem of \cref{sec:p17-partial-observation-decisions}.
Here the domain is the whole nonnegative cone, including
its origin, and the control may affect observation quality.

\section{Controlled observations and the nonnegative cone}

\begin{hbdefinition}{def:p17-astrom-ii-model}{The Part II model and its normalization}
Retain finite hidden states \(1,\ldots,n\), finite observations
\(1,\ldots,m\), a positive finite horizon \(N\), and a common
nonempty compact control set \(U\subseteq\mathbb R^d\).
At date \(t<N\), let \(P_t(u)=(p_{si,t}(u))\) and
\(Q_t(u)=(q_{ij,t}(u))\) be stochastic matrices.
The new dependence of \(Q_t\) on \(u\) means that the control
chosen after the current observation affects the observation
of the next hidden state. Its precise law is
\[
 P(X_{t+1}=i,Y_{t+1}=j\mid X_t=s,H_t,U_t=u)
                      =p_{si,t}(u)q_{ij,t}(u).
\]
The current control uses the observed history and earlier
controls, with the same private-randomization condition as
in \cref{def:p17-astrom-model}. Initial observation data,
if used, are specified separately before the first control.
Let \(g_t(u)\in\mathbb R^n\) be the immediate reward vector.
The entries of \(P_t,g_t\) are continuous in \(u\);
the entries of \(Q_t\) are Borel.
Continuity of \(Q_t\), when needed for attainment below,
is an additional explicit condition.
Put \(G_t=\sup_{u,i}|g_t(u)_i|<\infty\).

The source's domain and norm are
\[
 S=\mathbb R_+^n,\qquad \|w\|_1=\sum_i|w_i|,\qquad
 \Delta_n=\{w\in S:\|w\|_1=1\}.
\]
A probability belief belongs to \(\Delta_n\).
A general \(w\in S\) is an unnormalized nonnegative weight
vector; its value will be its mass times the value of its
normalized belief. Define the positive linear maps
\[
                (A_{j,t}(u)w)_i
                      =q_{ij,t}(u)\sum_s p_{si,t}(u)w_s .
\]
The source suppresses the \(t,u\) arguments of \(A_j,P,Q,g\).
For fixed \(t,u\), the map is linear in \(w\), even when
its dependence on the chosen control is nonlinear.
\end{hbdefinition}

\begin{hbdefinition}{def:p17-astrom-ii-recursions}{Finite maximum and minimum recursions}
Set \(V_{N+1}^+=V_{N+1}^-=0\).
The source uses \(V_t\) for either criterion; the superscripts
here identify its maximum and minimum presentations.
At the terminal date put
\[
              V_N^+(w)=\sup_{u\in U}\langle g_N(u),w\rangle,
 \qquad
              V_N^-(w)=\inf_{u\in U}\langle g_N(u),w\rangle.
\]
For \(t<N\), define
\[
 \begin{split}
 V_t^+(w)&=\sup_{u\in U}
       \left\{\langle g_t(u),w\rangle+
       \sum_{j:\|A_{j,t}(u)w\|_1>0}
       \|A_{j,t}(u)w\|_1
       V_{t+1}^+\!\left(\frac{A_{j,t}(u)w}{\|A_{j,t}(u)w\|_1}\right)
       \right\},\\
 V_t^-(w)&=\inf_{u\in U}
       \left\{\langle g_t(u),w\rangle+
       \sum_{j:\|A_{j,t}(u)w\|_1>0}
       \|A_{j,t}(u)w\|_1
       V_{t+1}^-\!\left(\frac{A_{j,t}(u)w}{\|A_{j,t}(u)w\|_1}\right)
       \right\}.
 \end{split}
\]
Set \(V_t^\pm(0)=0\); a zero observation column contributes
zero and is never divided by its norm.
Where the extrema are attained, these are exactly the
source's maximum equation (1.1) and minimum variant.
The supremum and infimum notation also states a definite
value when a discontinuous observation-control dependence
prevents attainment.
\end{hbdefinition}

\begin{hbtheorem}{thm:p17-astrom-ii-mass-bridge}{proved=annex:p17.astrom.ii.mass.bridge}{Mass conservation, homogeneous extension and the 1965 interface}
For every \(u\in U\), \(1\leq t<N\), and \(w\in S\),
\[
           \sum_j\|A_{j,t}(u)w\|_1=\|w\|_1.
\]
The recursions are finite valued and positively homogeneous:
\[
 |V_t^\pm(w)|\leq C_t\|w\|_1,\qquad
 V_t^\pm(cw)=cV_t^\pm(w)\quad(c\geq0),\qquad
 C_t=\sum_{k=t}^N G_k.
\]
Consequently they have the equivalent unnormalized form
\[
 V_t^\pm(w)=\mathop{\mathrm{ext}}_{u\in U}
           \left[\langle g_t(u),w\rangle+
                         \sum_jV_{t+1}^\pm(A_{j,t}(u)w)\right],
\]
where \(\mathrm{ext}\) is \(\sup\) for \(+\) and \(\inf\)
for \(-\), and the continuation is absent at \(N\).
On \(\Delta_n\) these are precisely the optimal finite-horizon
expected reward and cost values of the controlled-observation
model. For \(w\ne0\),
\[
              V_t^\pm(w)=\|w\|_1
                       V_t^\pm(w/\|w\|_1).
\]
When \(Q_t\) is independent of \(u\), the filter and recursion
reduce to the proved 1965 formulas.
\end{hbtheorem}

\section{The two source convexity lemmas}

\begin{hblemma}{thm:p17-astrom-ii-maxima}{proved=annex:p17.astrom.ii.maxima}{Lemma 1 and arbitrary control families}
For real numbers \(a,b,c,d\),
\[
             \max(a+b,c+d)\leq\max(a,c)+\max(b,d).
\]
The pointwise maximum of two convex real functions on
a convex set is convex.
More generally the pointwise supremum of any nonempty family
of convex functions, when finite valued, is convex.
Finite sums of convex functions are convex.
The corresponding statements for concavity use pointwise
infima and finite sums of concave functions.
The arbitrary-family version is needed when \(U\) is
an uncountable compact control set.
\end{hblemma}

\begin{hblemma}{thm:p17-astrom-ii-perspective}{proved=annex:p17.astrom.ii.perspective}{Lemma 2 on the closed cone, including vanishing columns}
Let \(A:S\to S\) be linear and let \(f\) be a finite
convex real function on \(\Delta_n\).
In particular this applies to the restriction of the
source's finite convex \(f:S\to\mathbb R\).
Define
\[
 ({\cal P}_A f)(w)=
 \begin{cases}
  \|Aw\|_1 f(Aw/\|Aw\|_1),&Aw\ne0,\\
  0,&Aw=0 .
 \end{cases}
\]
Then \({\cal P}_A f\) is convex and positively homogeneous
on all of \(S\). If \(f\) is concave, it is concave instead.
The definition and inequality include \(w=0\), zero columns
of \(A\), and mixture endpoints.
No division by a vanishing \(\|Aw\|_1\) occurs.
\end{hblemma}

\section{The main result and its minimum-cost counterpart}

\begin{hbtheorem}{thm:p17-astrom-ii-convexity}{proved=annex:p17.astrom.ii.convexity}{The Part II maximum-reward theorem}
For every \(1\leq t\leq N\), \(V_t^+\) is convex on
the complete nonnegative cone \(S\):
\[
 V_t^+(\theta w+(1-\theta)v)
          \leq\theta V_t^+(w)+(1-\theta)V_t^+(v)
                      \qquad(0\leq\theta\leq1).
\]
This includes the source's attained-maximum formulation.
Linearity and positivity of each fixed-control \(A_{j,t}(u)\),
the common control set, and finiteness of the value are the
essential algebraic premises. The conclusion does not require
linearity of \(P_t,Q_t\) as functions of the control.
More generally the same convexity conclusion holds for any
nonempty common control family, any positive linear maps
\(A_{j,t}(u):S\to S\), and any linear reward candidates,
provided all the recursively defined suprema are finite.
For this algebraic statement, probability normalization,
compactness and measurability are not needed; no policy
existence is inferred from that broader formulation.
\end{hbtheorem}

\begin{hbtheorem}{thm:p17-astrom-ii-concavity}{proved=annex:p17.astrom.ii.concavity}{The full minimum-cost result stated in the source's final remark}
For every \(1\leq t\leq N\), \(V_t^-\) is concave on \(S\):
\[
 V_t^-(\theta w+(1-\theta)v)
          \geq\theta V_t^-(w)+(1-\theta)V_t^-(v)
                      \qquad(0\leq\theta\leq1).
\]
Equivalently,
\(V_t^-[g]= -V_t^+[-g]\).
Thus the minimization in the 1965 switch example has the
concavity direction, while the production example's reward
maximization has the convexity direction.
The same concavity conclusion holds for the analogous
arbitrary common control families when all recursive infima
are finite.
\end{hbtheorem}

\begin{hbtheorem}{thm:p17-astrom-ii-continuity}{proved=annex:p17.astrom.ii.continuity}{Continuity on the full cone and attained control-dependent sensing}
For either sign, the value satisfies the explicit bound
\[
            |V_t^\pm(w)-V_t^\pm(v)|
                            \leq C_t\|w-v\|_1\qquad(w,v\in S).
\]
In particular its homogeneous extension is continuous at
the cone origin and all boundary faces.
If every entry of \(Q_t(u)\) is also continuous in \(u\),
every extremum is attained and a deterministic Borel
belief policy is optimal at every starting belief.
It can be chosen constant along positive rays of \(S\).
The Bayesian filter, policy equivalence, finite-horizon
verification and information comparisons continue to hold
with the declared control-dependent next-observation law.
Convexity alone does not establish an attained optimum when
that continuity condition is absent.
\end{hbtheorem}

\section{The source's examples and exact boundary calculations}

\begin{hbtheorem}{thm:p17-astrom-ii-examples}{proved=annex:p17.astrom.ii.examples}{The two original example references on the cone}
Part II refers to the examples of Part I and gives no new
numerical example. Their full calculations are in
\cref{thm:p17-astrom-example-one,thm:p17-astrom-example-two}.
For the switch example with \(q=4/5\), the stage-one
minimum cost on the cone is
\[
 V_1^-(w)=
 \min\{w_1,w_2,\tfrac9{25}w_1+\tfrac1{25}w_2,
                       \tfrac9{25}w_2+\tfrac1{25}w_1\}.
\]
It gives \(V_1^-(2,1)=11/25\).
For normalized beliefs \((1/10,9/10)\) and
\((9/10,1/10)\), the values are both \(9/125\);
their midpoint has value \(1/5\), exhibiting strict concavity.

For the production example, the stage-three maximum reward is
\[
                 V_3^+(w)=
                  \max\{\tfrac{55}{2}w_1,\
                                  \tfrac{49}{2}w_1+\tfrac92w_2\}.
\]
It gives \(V_3^+(2,1)=55\).
At the state point masses its values are \(27.5\) and \(4.5\);
their midpoint has value \(14.5<16\), exhibiting strict convexity.
Both cone values are zero at the origin and scale linearly
along every nonnegative ray.
\end{hbtheorem}

\begin{hbtheorem}{thm:p17-astrom-ii-boundaries}{proved=annex:p17.astrom.ii.boundaries}{Why the linear-map, norm and control premises cannot be discarded}
There are explicit examples showing each of the following.
\begin{enumerate}
\item A nonlinear map \(A:S\to S\) need not preserve the
      convexity conclusion of the source's theorem.
\item Replacing the additive \(\ell^1\) norm on \(S\) by
      the Euclidean norm can make the perspective of a
      convex function nonconvex.
\item A linear map not preserving the nonnegative cone
      can invalidate the perspective argument or leave
      its stated domain.
\item Allowing the feasible controls to depend on the
      belief can destroy convexity of the optimized value.
\item A compact control set and continuous state transitions
      and rewards do not ensure an attained optimum when
      the observation matrix depends discontinuously on control.
\end{enumerate}
The last example has a finite-horizon convex value and
arbitrarily accurate policies, but no optimal policy for
its fair initial belief.
\end{hbtheorem}

\begin{proofboundary}{def:p17-astrom-ii-boundary}{What Part II establishes about computation and limits}
The paper establishes finite-horizon convexity and concavity.
Its discussion of usefulness for numerical methods and convergence
does not give an infinite-horizon theorem, a numerical algorithm
with an error certificate, or a convergence theorem for a grid
approximation. The displayed recursion defines successive
optimization problems on the full information-state domain.
Effective evaluation for an arbitrary continuous-control input
also requires a specified representation, optimization method
and approximation guarantees.
General piecewise-linear representations and later approximation
results retain their own sources and proofs.
\end{proofboundary}
